\chapter{คลังชุดคำสั่งวิศวกรรมพรอมพต์สำหรับการพัฒนาแอปพลิเคชัน}
\label{app:prompts}

ภาคผนวกนี้รวบรวมคลังชุดคำสั่งวิศวกรรมพรอมพต์ (Prompt Engineering Corpus) ตามกรอบการทำงาน CLEAR Prompting Framework เพื่อให้ผู้เริ่มต้น นักศึกษา และนักพัฒนา สามารถนำไปใช้สั่งการโมเดลปัญญาประดิษฐ์เชิงกำเนิด (เช่น Claude 3.7 Sonnet, Gemini 1.5/2.0 Flash/Pro, หรือ GPT-4o) ในการสร้างโค้ดและแก้ไขปัญหาในแต่ละโมดูลได้อย่างถูกต้อง แม่นยำ และเป็นไปตามหลักการทางวิทยาศาสตร์

\section{ชุดพรอมพต์ออกแบบสถาปัตยกรรมระบบ Clean Architecture}

\begin{promptbox}[title={พรอมพต์ ก.1: การวางโครงสร้างโปรเจกต์ Flutter Clean Architecture}]
\begin{lstlisting}
คุณคือสถาปนิกซอฟต์แวร์อาวุโสด้าน Mobile Embedded System
จงออกแบบโครงสร้างโปรเจกต์ Flutter สำหรับแอปพลิเคชันวิเคราะห์ดิน JC Digital Soil AI 
โดยใช้หลักการ Clean Architecture ที่แบ่งความรับผิดชอบอย่างชัดเจน:
1. Core Layer: ค่าคงที่, ธีมสี Cyber Neon Dark Glow, เกณฑ์วิเคราะห์ดินทางปฐพีวิทยา
2. Data Layer:
   - บริการสื่อสารพอร์ต USB Type-C OTG Serial (CH340/MAX485)
   - โมดูลถอดรหัสเฟรมข้อมูล Modbus RTU 8-in-1 และตรวจสอบ CRC-16
   - บริการบันทึกและส่งออกข้อมูลประวัติลงไฟล์ Excel (.xlsx)
3. Domain Layer:
   - เอนจินโมเดลโครงข่ายประสาทเทียมอิงฟิสิกส์ (Physics-Informed Neural Network: PINN)
   - โมเดลฟิสิกส์การปรับเทียบอุณหภูมิ (Arrhenius EC) และความชันเนิร์นสท์ (Nernstian pH)
   - เอนจินวิเคราะห์ความอุดมสมบูรณ์ของดินและให้คำแนะนำสำหรับสวนทุเรียน
4. Presentation (UI) Layer:
   - แดชบอร์ดแสดงผลค่าพารามิเตอร์ 8 ชนิดแบบเรียลไทม์
   - วิดเจ็ตมาตรวัดวงกลมเรืองแสง (Glassmorphic Circular Gauge)
   - หน้าต่างการตั้งค่าพอร์ตสื่อสารและคู่มือการเชื่อมต่อ OTG
จงแสดงโครงสร้างไดเรกทอรีไฟล์ทั้งหมดในรูปแบบ Tree พร้อมคำอธิบายหน้าที่ของแต่ละไฟล์
\end{lstlisting}
\end{promptbox}

\section{ชุดพรอมพต์สร้างไดรเวอร์ USB OTG และถอดรหัส Modbus RTU}

\begin{promptbox}[title={พรอมพต์ ก.2: การเขียนโมดูลถอดรหัส Modbus RTU และตรวจสอบ CRC-16}]
\begin{lstlisting}
คุณคือวิศวกรระบบสมองกลฝังตัวผู้เชี่ยวชาญโปรโตคอลอุตสาหกรรม
จงเขียนคลาส ModbusParser และ ModbusFrame ภาษา Dart สำหรับถอดรหัสข้อมูลจากหัววัดดิน 8-in-1:
1. รองรับโครงสร้างเฟรมมาตรฐาน Modbus RTU:
   [Slave ID (1B), Function Code (1B), Byte Count (1B), Data (14B หรือ 16B), CRC-16 (2B)]
2. เขียนอัลกอริทึมคำนวณและตรวจสอบ CRC-16 Modbus (พหุนาม 0xA001) แบบบิตต่อบิต
3. ใช้ ByteData เพื่ออ่านค่ารีจิสเตอร์ขนาด 16-bit Unsigned Integer (Big-Endian)
4. แปลงค่าตามตัวคูณมาตราส่วน (Scale Factors) ของหัววัด:
   - อุณหภูมิดิน: หารด้วย 10.0 (องศาเซลเซียส)
   - ความชื้นดิน: หารด้วย 10.0 (เปอร์เซ็นต์)
   - สภาพนำไฟฟ้า: ค่าจำนวนเต็ม (ไมโครซีเมนต์ต่อเซนติเมตร)
   - ความเป็นกรด-ด่าง: หารด้วย 10.0 หรือ 100.0 (pH)
   - ไนโตรเจน, ฟอสฟอรัส, โพแทสเซียม: ค่าจำนวนเต็ม (มิลลิกรัมต่อกิโลกรัม)
5. หากตรวจสอบ CRC-16 ไม่ผ่าน ให้โยนข้อยกเว้น FormatException พร้อมแสดง Hex Dump
\end{lstlisting}
\end{promptbox}

\section{ชุดพรอมพต์สร้างโมเดลโครงข่ายประสาทเทียมอิงฟิสิกส์ (PINN)}

\begin{promptbox}[title={พรอมพต์ ก.3: การสร้างโมเดล Two-Stage Decoupling PINN บนอุปกรณ์ฝังตัว}]
\begin{lstlisting}
คุณคือนักวิทยาศาสตร์ข้อมูลด้านฟิสิกส์เกษตรดิจิทัล (Digital Agriphysics Data Scientist)
จงเขียนคลาส PinnSoilEngine ภาษา Dart เพื่อชดเชยค่าความคลาดเคลื่อนของหัววัดดิน:
1. การชดเชยขั้นตอนที่ 1 (First-Principles Physics Stage):
   - ชดเชยการลอยเลื่อนของค่า EC ตามอุณหภูมิอ้างอิง 25 องศาเซลเซียส ด้วยสมการอาร์เรเนียส
   - ชดเชยค่า pH ตามความชันของเนิร์นสท์ S_N(T) = 2.3026 * R * T / F
   - คำนวณสภาพยอมทางไฟฟ้าสัมพัทธ์ของดินจากสมการเชิงประจักษ์ของท็อปป์ (Topp's Model)
2. การชดเชยขั้นตอนที่ 2 (Residual MLP PINN Stage):
   - โครงข่ายประสาทเทียมขนาดเล็ก (TinyML MLP) ชดเชยความไม่เป็นเชิงเส้นของสารละลายดิน
   - ใช้ฟังก์ชันกระตุ้น LeakyReLU เพื่อลดปัญหา Vanishing Gradient
   - ควบคุมขอบเขตผลลัพธ์ด้วยคำสั่ง clamp เพื่อป้องกันค่าหลุดออกนอกช่วงทางกายภาพ
3. ฟังก์ชันต้องคืนค่า Map<String, double> ประกอบด้วยค่าหลังปรับเทียบและค่าส่วนต่าง (Delta)
\end{lstlisting}
\end{promptbox}

\section{ชุดพรอมพต์ออกแบบหน้าจอแดชบอร์ดและวิดเจ็ตมาตรวัด}

\begin{promptbox}[title={พรอมพต์ ก.4: การสร้างวิดเจ็ตมาตรวัดวงกลม CustomPainter สไตล์โมเดิร์น}]
\begin{lstlisting}
คุณคือผู้เชี่ยวชาญการออกแบบ Flutter UI/UX สไตล์ Glassmorphism และ Cyberpunk Theme
จงสร้างคลาส SoilGaugeCard และ CustomPainter สำหรับแสดงผลพารามิเตอร์ดินแต่ละตัว:
1. มีวงแหวนความคืบหน้า (Circular Arc) 240 องศา แสดงระดับค่าปัจจุบันเทียบกับช่วงปกติ
2. ปรับเปลี่ยนสีเส้นโค้งเรืองแสงอัตโนมัติตามเกณฑ์ทางการเกษตร:
   - ต่ำกว่าเกณฑ์ (Under-optimal): สีส้มอำพัน
   - เหมาะสม (Optimal): สีเขียวมรกตนีออน
   - เกินเกณฑ์ (Over-optimal): สีแดงสดแจ้งเตือน
3. มีการ์ดแบบ Frosted Glass พื้นหลังสีกึ่งโปร่งแสง พร้อมขอบเส้นนีออนบางเบา
4. แสดงตัวเลขขนาดใหญ่ชัดเจน พร้อมหน่วยวัด และข้อความสถานะภาษาไทย (เช่น "ปกติ", "ดินกรดจัด")
5. รองรับการกดแตะเพื่อเปิดหน้าต่างแสดงกราฟแนวโน้มและคำแนะนำทางปฐพีวิทยา
\end{lstlisting}
\end{promptbox}

\section{ชุดพรอมพต์แก้ไขปัญหาการเชื่อมต่อ RS485 Half-Duplex บนฮาร์ดแวร์จริง}

\begin{promptbox}[title={พรอมพต์ ก.5: การแก้ปัญหาสายแปลง USB-RS485 มิวต์ตัวรับสัญญาณ (RTS DE/RE)}]
\begin{lstlisting}
คุณคือผู้เชี่ยวชาญด้านฮาร์ดแวร์ไดรเวอร์ USB Serial บนระบบปฏิบัติการแอนดรอยด์
บริบท: แอปพลิเคชันเชื่อมต่อสายแปลง USB-to-RS485 (CH340/MAX485) ผ่าน USB OTG สำเร็จ
แต่ไม่ได้รับข้อมูลตอบกลับจากหัววัดดิน (Rx = 0 Bytes ตลอดเวลา)
จงวิเคราะห์สาเหตุและเขียนโค้ดภาษา Dart เพื่อแก้ไขปัญหา:
1. อธิบายกลไกการทำงานของสายสัญญาณ RTS (Request to Send) ที่เชื่อมกับขา DE/~RE ของชิป MAX485
2. แก้ไขโค้ดการกำหนดค่าพอร์ต UsbPort:
   - ต้องสั่ง port.setRTS(false) เพื่อปลดให้ขา RE เป็น Low เปิดรับสัญญาณจากหัววัด
   - ต้องกำหนด Flow Control เป็น FLOW_CONTROL_OFF
3. เพิ่มระบบ Auto-Baud Rate ตรวจจับและสลับระหว่าง 4800 bps และ 9600 bps อัตโนมัติ
4. เพิ่มการแสดงผลข้อมูลทางเทคนิคสด (Telemetry): จำนวนแพ็กเก็ต Tx, ไบต์ Rx, และรหัส Hex ล่าสุด
\end{lstlisting}
\end{promptbox}
